\section{A quantum-versus-classical resilience crossover map}\label{sec:crossover}

The contribution of this paper is not another single-point claim that a quantum
sensor outperforms its classical counterpart. It is a \emph{map}: for each
quantum lever, a grid of operating points across the realistic published-parameter
envelope, at every node of which we report the resilience figure of merit as a
distribution rather than a number, and against which we locate the
\emph{break-even} where the quantum advantage disappears. The question the map
answers is the one the field's own reviews say is open
\citep{quantumpntreview2026}: not \emph{whether} a cold-atom inertial sensor or an
optical clock can help during a GNSS outage, but \emph{when} --- under what outage
duration, platform environment, and device maturity --- it actually changes the
outcome, and when a same-maturity classical unit is the better choice.

Every result in this section is produced by the same machinery the validated packs
use. For each lever we sweep the operating grid with \code{nd\_sweep\_ensemble};
each node is a Monte-Carlo ensemble (64 seeds for the inertial study, 32 for the
clock) over the parameter envelope, and the figure of merit is summarised by its
mean, percentiles, and a 95\% bootstrap confidence interval computed by the same
bootstrap the ensemble packs apply elsewhere. The quantum and classical legs share
one neutral code path, swapped by parameters alone, so the comparison is never made
across two implementations. No new physics is introduced. The cold-atom noise is a
first-principles interferometer model --- a shot-noise term $1/(C\sqrt{N})$ in atom
number $N$ and contrast $C$, combined in quadrature with a platform-vibration floor
$\sqrt{S_a/(3T)}$ in the acceleration PSD $S_a$ and interrogation time $T$ --- and
the classical leg uses the NaveGo-validated inertial error profile
\citep{gonzalez2017navego}. We restate the binding caveat of
Section~\ref{sec:limitations} once here and rely on it throughout: \textbf{every
quantum figure in this section is a performance-model projection from
laboratory-maturity device coefficients, not a measurement.}

\subsection{Inertial: cold-atom versus navigation-grade dead reckoning}

The inertial lever pits a cold-atom accelerometer of the Exail/Templier class
\citep{templier2022} against a navigation-grade strapdown unit
\citep{groves2013}, both dead-reckoning position through a GNSS outage. We sweep
two axes that dominate the outcome: outage duration, from \SI{30}{\second} to
\SI{1800}{\second}, and the platform-vibration acceleration PSD, from
$1~\mu\mathrm{g}/\sqrt{\mathrm{Hz}}$ (an isolated, quiescent platform) to roughly
$3~\mathrm{mg}/\sqrt{\mathrm{Hz}}$ (a noisy host). The figure of merit is the
95th-percentile horizontal position error over the ensemble; the advantage at a
node is the classical 95th-percentile error divided by the quantum one.
Figure~\ref{fig:inertial} shows the full map.

\begin{figure}[t]
  \centering
  \includegraphics[width=0.95\linewidth]{inertial.pdf}
  \caption{Inertial crossover map. Cold-atom accelerometer (Exail/Templier class,
  \citealp{templier2022}) versus navigation-grade strapdown
  (\citealp{groves2013}), dead-reckoning 95th-percentile position error as a
  function of outage duration and platform-vibration acceleration PSD. Each node is
  a 64-seed Monte-Carlo ensemble with a 95\% bootstrap confidence interval; the
  dashed contour is the break-even (advantage~$=1$) where the navigation-grade unit
  matches the cold-atom sensor. The advantage is overwhelming on a quiet platform
  and erodes with vibration; the break-even PSD rises with outage duration. The
  quantum leg is a performance-model projection from laboratory-maturity device
  coefficients, not a measurement.}
  \label{fig:inertial}
\end{figure}

The first finding is that on an isolated platform the cold-atom advantage is large
and grows with outage length. At $1~\mu\mathrm{g}/\sqrt{\mathrm{Hz}}$ the advantage
is a factor of $307$ for a \SI{30}{\second} outage (cold-atom
\SI{2.5}{\milli\meter} versus navigation-grade \SI{0.75}{\meter} at the 95th
percentile), rising to $439$ at \SI{60}{\second} and to $1915$ at
\SI{1800}{\second}, where the navigation-grade error has grown to
\SI{2.3}{\kilo\meter} while the cold-atom error is near \SI{1.2}{\meter}. The
confidence intervals are tight relative to this gap --- the cold-atom 95th
percentile at the \SI{30}{\second}, $1~\mu\mathrm{g}/\sqrt{\mathrm{Hz}}$ node has a
bootstrap interval of $[2.0,\,2.9]~\si{\milli\meter}$ --- so the advantage on a
quiet platform is not a sampling artefact.

The second finding is that this advantage is conditional, and the map shows exactly
where it fails. As the vibration PSD rises the cold-atom error grows steeply, driven
by the $\sqrt{S_a/(3T)}$ vibration floor, while the navigation-grade error is
comparatively insensitive to it. The two curves cross: for a short
\SI{30}{\second} outage at a vibration level of
$\approx 405~\mu\mathrm{g}/\sqrt{\mathrm{Hz}}$ (about $0.4~\mathrm{mg}/\sqrt{\mathrm{Hz}}$)
the advantage falls to $0.75$ --- below unity, meaning the navigation-grade unit
\emph{wins}. On a vibrating host during a brief outage, the classical sensor is the
better choice, and the map says so plainly rather than reporting only the favourable
corner.

The third finding ties the first two together: the break-even vibration level rises
\emph{monotonically} with outage duration, from
$334~\mu\mathrm{g}/\sqrt{\mathrm{Hz}}$ at \SI{30}{\second} to
$2640~\mu\mathrm{g}/\sqrt{\mathrm{Hz}}$ at \SI{1800}{\second} (Table~\ref{tab:breakeven}).
The longer the outage, the more vibration the cold-atom sensor can tolerate before
losing. The mechanism is in the model: the cold-atom advantage over the
navigation-grade unit is dominated by the latter's accumulating bias error, which
grows roughly as $t^2$, whereas the cold-atom sensor's own vibration-limited error
is a random walk growing roughly as $t^{1.5}$. The bias term outruns the random-walk
term, so the duration at which the quantum sensor pays off increases the more
vibration it must endure --- the longer the outage, the wider the environmental
window in which it is the right tool.

\begin{table}[t]
  \centering
  \caption{Break-even platform-vibration PSD at which the navigation-grade unit
  matches the cold-atom sensor, as a function of outage duration. The break-even
  rises monotonically: a longer outage tolerates more vibration before the quantum
  advantage is lost.}
  \label{tab:breakeven}
  \begin{tabular}{rr}
    \toprule
    Outage duration (s) & Break-even PSD ($\mu\mathrm{g}/\sqrt{\mathrm{Hz}}$) \\
    \midrule
    30   & 334  \\
    60   & 448  \\
    120  & 724  \\
    300  & 1168 \\
    600  & 1565 \\
    1200 & 2250 \\
    1800 & 2640 \\
    \bottomrule
  \end{tabular}
\end{table}

\subsection{Clock: a same-maturity holdover ladder}

The clock lever asks how long each of three clock technologies can hold timing to a
\SI{1}{\micro\second} budget after its disciplining GNSS reference is lost. We
deliberately frame this as a technology-readiness ladder rather than a single
contest, because the three rungs are not the same maturity and a like-for-like
winner-flip does not exist here: a ground-laboratory strontium optical-lattice clock
\citep{origlia2015}, a flight-qualified active hydrogen maser of the ACES/PHARAO
class (the relevant point being that this class has actually flown, on the ISS), and
a deployed chip-scale atomic clock (CSAC). The figure of merit is the
95th-percentile holdover timing error against the \SI{1}{\micro\second} threshold;
each rung is a 32-seed ensemble with a bootstrap interval.
Figure~\ref{fig:clock} shows the three curves.

\begin{figure}[t]
  \centering
  \includegraphics[width=0.8\linewidth]{clock.pdf}
  \caption{Clock holdover ladder. 95th-percentile timing error against a
  \SI{1}{\micro\second} budget as a function of holdover duration, for a
  ground-laboratory strontium optical-lattice clock (\citealp{origlia2015}), a
  flight-qualified ACES/PHARAO-class hydrogen maser, and a deployed chip-scale
  atomic clock. Each point is a 32-seed ensemble with a 95\% bootstrap interval and
  carries an inline technology-readiness tag. There is no winner-flip: the optical
  clock dominates stability, but its margin is a laboratory figure while the maser is
  the realistically flown option. The CSAC reaches the budget at $\approx$\SI{8.5}{\hour}.}
  \label{fig:clock}
\end{figure}

There is no crossover here, and reporting that honestly is the point. The
optical-lattice clock dominates the stability ranking at every holdover duration,
and both it and the maser hold the \SI{1}{\micro\second} budget essentially
indefinitely over the swept range: at \SI{24}{\hour} the optical clock's
95th-percentile holdover error is sub-nanosecond and the maser's is well under a
nanosecond, both three orders of magnitude inside budget. But the honest comparison
is the maturity-aware one. The optical clock's margin is a \emph{ground-laboratory}
figure; the maser is the option that has actually flown, and is therefore the
realistic on-orbit choice today. The deployed CSAC, by contrast, reaches the
\SI{1}{\micro\second} budget at $\approx$\SI{8.5}{\hour} of holdover --- a result
that matters operationally precisely because the CSAC is the rung that is already
fielded. The value of the ladder is that it refuses to compare a laboratory optical
clock against a deployed CSAC as if they were interchangeable, and instead pairs each
figure with the maturity tag that determines whether it is available to a mission.

\subsection{Time transfer: optical versus RF with a non-reciprocity budget}

The third lever compares optical inter-satellite time transfer against RF two-way
satellite time-and-frequency transfer (TWSTFT), and the methodological care is in
refusing the easy comparison. A jitter-only optical floor would report a figure that
no real two-way link achieves, because the limiting error in operational time
transfer is not jitter but \emph{non-reciprocity}: equipment-delay asymmetry between
the two terminals and path variation that does not cancel in the two-way difference.
We therefore model a non-reciprocal differential-delay budget for the optical link
\citep{giorgetta2013,deschenes2016} and adopt a conservative \SI{1}{\pico\second}
on-orbit optical target rather than the \SI{1}{\femto\second} laboratory floor, so
that the optical figure is charged with the same class of systematic that limits the
RF figure. Under this honest accounting the optical link reaches a range-equivalent
RMS of $\approx$\SI{0.29}{\milli\meter} (non-reciprocity included, a
\SI{1}{\pico\second} on-orbit projection target rather than a measured link budget)
against $\approx$\SI{341}{\milli\meter} for the RF link --- still a
three-order-of-magnitude advantage, but one obtained from a like-for-like budget
rather than from comparing an optical jitter floor against an RF link burdened with
systematics.

\subsection{What the map establishes}

Taken together, the three levers answer the question of \emph{when} quantum sensing
changes a PNT-resilience outcome, rather than asserting that it always does. The
inertial map shows a large but conditional advantage with a vibration-dependent
break-even that the classical sensor wins inside; the clock ladder shows a clear
stability ordering that is nonetheless governed by maturity, with no winner-flip to
overstate; and the time-transfer comparison shows a robust optical advantage only
once the optical figure is charged with the non-reciprocal systematics that limit
real links. Each figure in this section is reproducible from a committed scenario and
fixed seed using the commands listed in Appendix~A; the inertial and clock grids,
their break-even values, and their confidence intervals regenerate bit-for-bit from
the same machinery used for the validated packs.
